\textbf{Verificação dimensional:}
\[
[P]=\mathrm{W}=\mathrm{J\,s^{-1}},\qquad \left[\frac{dE_k}{dt}\right]=\mathrm{J\,s^{-1}}.\quad\checkmark
\]

\textbf{Dedução:}

A potência mecânica instantânea de uma força $\mathbf{F}$ agindo sobre um corpo em movimento com velocidade $\mathbf{v}$ é definida como:
\[
P = \mathbf{F}\cdot\mathbf{v} = Fv\cos\theta,
\]
onde $\theta$ é o ângulo entre força e velocidade. Para movimento unidimensional:
\[
P = Fv.
\]

A energia cinética é:
\[
E_k=\frac{1}{2}mv^2.
\]

Derivando em relação ao tempo, pela regra da cadeia:
\[
\frac{dE_k}{dt}=\frac{d}{dt}\!\left(\frac{1}{2}mv^2\right)=\frac{1}{2}m\cdot 2v\frac{dv}{dt}=mv\frac{dv}{dt}.
\]

Pela Segunda Lei de Newton, $F=ma=m\dfrac{dv}{dt}$, logo:
\[
\frac{dE_k}{dt}=mv\frac{dv}{dt}=\underbrace{m\frac{dv}{dt}}_{=F}\cdot v = Fv = P.\quad\blacksquare
\]

Portanto, a potência é a taxa temporal de variação da energia cinética, $P=\dfrac{dE_k}{dt}$.

